% ============================================================================
%  【已废弃 · 已合并】 DEPRECATED -- MERGED INTO reconstruct
%  本稿已与 microfusion.tex（融合+蒸馏篇）合并重构，取代稿见同目录：
%      reconstruct.tex     (英文, 期刊长文)
%      reconstruct_zh.tex  (中文)
%  本文件仅作合并前的历史留档，不再维护。特别地：本稿摘要/引言/结论中
%  「gate-style 最优」与「语料规模边界成立」的表述已被 v44/v45 审计推翻，
%  合并稿中已按「可容许性 vs 实证」两层结论改写。请勿引用本文件。
%  合并日期：2026-10-11
% ============================================================================
\documentclass[10pt,twoside]{article}

% ==================== NeurIPS 官方样式 ====================
\usepackage{neurips_2025}

% ==================== 数学与符号 ====================
\usepackage{amsmath,amssymb,amsfonts,amsthm}
\usepackage{mathtools}
\usepackage{bm}
\usepackage{siunitx}

% ==================== 图表 ====================
\usepackage{graphicx}
\usepackage{xcolor}
\usepackage{booktabs}
\usepackage{multirow}
\usepackage{array}
\usepackage{colortbl}
\usepackage{tikz}
\usetikzlibrary{arrows.meta,positioning,calc,fit,shapes.geometric}
\usepackage{caption}
\usepackage{subcaption}
\usepackage{wrapfig}

% ==================== 链接与引用 ====================
\usepackage[colorlinks=true, linkcolor=black, citecolor=black, urlcolor=black]{hyperref}
\usepackage{url}
\usepackage{doi}
\usepackage{enumitem}
\usepackage{footnote}

% ==================== 自定义命令 ====================
\newcommand{\todo}[1]{\textcolor{red}{[TODO: #1]}}
\newcommand{\fixme}[1]{\textcolor{orange}{[FIXME: #1]}}

% 数学符号
\newcommand{\E}{\mathbb{E}}
\newcommand{\R}{\mathbb{R}}
\newcommand{\argmin}{\operatorname{argmin}}
\newcommand{\argmax}{\operatorname{argmax}}
\newcommand{\softmax}{\operatorname{softmax}}
\providecommand{\keywords}[1]{\vspace{0.5em}\noindent\textit{Keywords:} #1}
\newcommand{\softplus}{\operatorname{softplus}}
\newcommand{\sigmoid}{\operatorname{sigmoid}}

% 版本引用
\newcommand{\vxx}[1]{\texttt{v#1}}

% ==================== 定理环境 ====================
\theoremstyle{definition}
\newtheorem{definition}{Definition}[section]
\newtheorem{assumption}{Assumption}[section]
\newtheorem{theorem}{Theorem}[section]
\newtheorem{lemma}{Lemma}[section]
\newtheorem{corollary}{Corollary}[section]
\newtheorem{proposition}{Proposition}[section]
\newtheorem{remark}{Remark}[section]
\newtheorem{example}{Example}[section]

% ==================== 元信息 ====================
\title{Temperature-Modulated Routing via Gate-Style Central Coordination in Heterogeneous Model Fusion}

\author{
    Anonymous Author(s)\\
    Affiliation\\
    \texttt{email@example.com}
}

\begin{document}

\maketitle

\vspace{-0.5em}
\begin{center}
\fbox{\parbox{0.94\textwidth}{\centering\small\textbf{DEPRECATED --- MERGED.} This draft has been merged and restructured into \texttt{reconstruct.tex} (English) and \texttt{reconstruct\_zh.tex} (Chinese), in this directory. It is kept for archival only: its numbers, section structure, and claims are superseded by the merged manuscript---in particular, the ``gate-style is optimal'' and ``corpus-scale boundary'' statements in the abstract, introduction, and conclusion below were overturned by the later held-out audits and are rewritten in the merged paper. Do not cite this file.}}
\end{center}
\vspace{0.5em}

\begin{abstract}
Mixture-of-experts (MoE) routing dynamically selects specialized models, but it remains unclear whether and how a \textbf{global hub} should coordinate expert selection. We identify a fundamental limitation: adding a learnable central vector to router logits destroys the relative ordering of already-converged logits, leading to performance collapse. To address this, we introduce \textbf{Gate-Style Central Coordination}, where the hub controls the softmax temperature rather than the logits themselves. This yields a \textbf{perfect degeneration} guarantee---disabling the hub incurs zero overhead---and provably maintains logit ordering. Across 29 engineering versions and 175+ experiments, we show that gate-style uniquely satisfies five independent design constraints that other hub modes violate. We diagnose a hidden failure mode---tanh saturation in long training severs the gradient chain to the hub---and establish a \textbf{corpus-scale boundary} showing that the optimal hub mode switches from adaptive-EMA to gate-style as data scale grows.
\end{abstract}

\keywords{Heterogeneous Model Fusion, Mixture of Experts, Temperature Scaling, Central Coordination, Gradient Flow}

\section{Introduction}

Mixture-of-experts (MoE) routing dynamically selects from a pool of specialized models \citep{shazeerOutrageouslyLargeNeural2017}. In heterogeneous model fusion, experts differ in architecture, pretraining objective, and modality. While routing mechanisms can select experts per token, they operate \textit{locally}: each token independently chooses which experts to activate, with no global coordination across the ensemble.

A natural question arises: \textbf{should a global hub coordinate expert selection?} Intuitively, a hub could broadcast summary information, resolve conflicts between experts, or modulate routing confidence. However, as we demonstrate, the simplest hub design---adding a learnable central vector directly to router logits---fails catastrophically. The root cause is not the hub concept itself, but \textit{how} the hub interacts with the routing mechanism. Direct additive bias perturbs the relative ordering of already-converged logits, destroying routing decisions and causing performance collapse.

We propose \textbf{Gate-Style Central Coordination}: instead of adding a bias to logits, the hub controls the \textit{temperature} of the softmax. This guarantees a \textbf{perfect degeneration} property---when disabled, the system is exactly equivalent to the no-hub baseline---and provably preserves logit ordering. We formalize five design constraints and prove gate-style is the only mode satisfying all five simultaneously.

Through systematic experimentation across 29 engineering versions and 175+ end-to-end experiments, we make the following contributions:

\begin{enumerate}
    \item \textbf{Mathematical Theory}: We formalize Gate-Style Central Coordination, prove perfect degeneration and order preservation theorems, and establish five design constraints (G1--G5) that uniquely characterize the gate-style mode.
    \item \textbf{Root Cause Diagnosis}: We identify a hidden failure mode in long-horizon training: tanh saturation in the internal temperature scalar severs the gradient chain to the hub, causing central parameters to stop updating.
    \item \textbf{Repair Schemes}: We propose and validate two repair directions---Softplus and Gradient-Scale---that restore gradient flow while preserving the perfect degeneration property.
    \item \textbf{Practical Guidelines}: We establish a \textbf{corpus-scale boundary} showing that gate-style is optimal for medium-scale data ($\geq 32$ training samples), while adaptive-EMA dominates under extreme scale constraints.
\end{enumerate}



\section{Related Work}

Mixture-of-Experts (MoE) layers dynamically route tokens to specialized feed-forward networks \citep{shazeerOutrageouslyLargeNeural2017}. In heterogeneous settings, experts differ in architecture or modality \citep{vaswani2017}, but most routing mechanisms operate \textit{locally}: each token independently selects experts with no global coordination.

Temperature scaling controls the sharpness of softmax distributions \citep{hintonDistillingKnowledgeNeural2015}. We are the first to use temperature modulation as a \textit{coordination mechanism} for a global hub in heterogeneous MoE routing.

Our work is inspired by Global Workspace Theory \citep{baarsCognitiveTheoryConsciousness1988} and centralized gating mechanisms \citep{wangLinformerSelfAttentionLinear2020}, but we focus on a key design choice: whether the hub should \textit{add} to router logits (broadcast) or \textit{scale} them (gating). We prove that additive bias can destroy logit ordering, while multiplicative modulation via temperature scaling preserves it by construction.

At the system level, Mixture-of-Agents \citep{wang2024moa} demonstrates that layered collaboration among heterogeneous agents, with an aggregator synthesizing their collective outputs---the modern form of the stacked-generalization idea \citep{wolpertStackedGeneralization1992}---is itself a source of capability. Our central-coordination question carries this agent-level insight into sub-layer routing: how should a hub coordinate experts that differ in architecture and modality?

Along a different axis, linear attention \citep{katharopoulos2020} and state-space models \citep{gu2022s4,gu2023mamba} replace softmax attention with fixed-size recurrent states that mix information across \textit{sequence positions}. These global mixers are orthogonal to our hub along two dimensions. First, the mixing axis: they compress the \textit{temporal} context of one model, whereas a central hub coordinates across the \textit{expert} axis of many frozen models at one position. Second, the interaction mode: global mixers aggregate \textit{content} into a carried state, while gate-style coordination modulates routing \textit{certainty}---a scalar temperature that rescales logits without adding content---which preserves logit ordering by construction and degenerates exactly when disabled. Content-carrying states satisfy neither property in general: a state accumulated across steps and added to logits can reorder expert preferences arbitrarily, and it cannot be removed without changing the model. The two mechanisms are therefore composable rather than competing; whether a temperature-modulated hub can coordinate a pool of linear-attention or state-space experts is an open direction.

\section{Problem Formulation}

Consider a heterogeneous model fusion system with $M$ experts. At each layer, the routing mechanism computes logits $z_m = w_m^\top x$ and weights $\alpha_m = \text{softmax}(z)_m$. The fused output is $y = \sum_m \alpha_m f_m(x)$.

We ask: should a global hub $c \in \mathbb{R}^{D_{\text{shared}}}$ influence routing, and how? We formalize five design constraints that any hub mechanism should satisfy:

\begin{itemize}
    \item \textbf{(G1) Perfect Degeneration}: $c=0$ yields the no-hub baseline exactly.
    \item \textbf{(G2) Order Preservation}: Hub must preserve relative logit ordering.
    \item \textbf{(G3) Monotonicity}: Hub intensity changes routing determinism monotonically.
    \item \textbf{(G4) Boundedness}: Hub influence has a theoretical upper bound.
    \item \textbf{(G5) Low Overhead}: Hub parameters are $O(D_{\text{shared}})$, negligible compared to experts.
\end{itemize}

We prove in Section~\ref{sec:theory} that only gate-style satisfies all five simultaneously.

\section{Gate-Style Central Coordination}

\subsection{Four Hub Modes}
\label{sec:modes}

We consider four ways a global hub can influence routing logits:

\begin{definition}[Broadcast Mode]
    $z_m^{(bc)} = z_m + (U c)_m$, where $U \in \mathbb{R}^{M \times D_{\text{shared}}}$ is learnable. The hub directly adds a bias to each expert's logit.
\end{definition}

\begin{definition}[Router-Norm Mode]
    $z_m^{(rn)} = \text{LayerNorm}(z_m + U c)$. Broadcast then normalize.
\end{definition}

\begin{definition}[Adaptive-EMA Mode]
    $\text{decay} = \sigma(\alpha \cdot \|c\|)$. The hub controls its own EMA decay rate.
\end{definition}

\begin{definition}[Gate-Style Mode (Proposed)]
    $t_{\text{scalar}} = W_t^\top c$, $\tau = 1 + \epsilon \cdot \tanh(t_{\text{scalar}})$, $z_m^{(gate)} = z_m / \tau$. The hub modulates softmax temperature.
\end{definition}

\subsection{Core Theorems}
\label{sec:theory}

\begin{theorem}[Perfect Degeneration]
    When $c = 0$, gate-style is exactly equivalent to the no-hub baseline.
\end{theorem}

\begin{proof}[Proof sketch]
    $c = 0 \Rightarrow t_{\text{scalar}} = 0 \Rightarrow \tau = 1$. Thus $z_m^{(gate)} = z_m / 1 = z_m$. Full proof in Appendix~\ref{app:proofs}.
\end{proof}

\begin{theorem}[Order Preservation]
    Gate-style preserves logit ordering; broadcast may reverse it.
\end{theorem}

\begin{proof}[Proof sketch]
    If $z_1 > z_2$, broadcast gives $z_1+U_1c$ vs $z_2+U_2c$, which can reverse if $|U_1c-U_2c| > z_1-z_2$. Gate-style divides both by $\tau>0$, preserving order. Full proof in Appendix~\ref{app:proofs}.
\end{proof}

\begin{theorem}[Unification]
    Among the four hub modes of \S\ref{sec:modes}, gate-style is the only one satisfying all five constraints (G1)--(G5).
\end{theorem}

\subsection{Five Hypotheses}

We evaluate five root-cause hypotheses for broadcast failure. Gate-style resolves all five: (H1) bypasses $U$ training dynamics; (H2) has no EMA delay via direct $W_t$ participation; (H3) guarantees baseline-equivalent behavior early via perfect degeneration; (H4) has $O(D_{\text{shared}})$ overhead; (H5) naturally determines modulation position via temperature.

\section{Diagnosis: Gradient Vanishing in Long Training}
\label{sec:diagnosis}

Systematic experiments across 29 engineering versions reveal a hidden failure mode: under long-horizon training, the central hub stops updating. Diagnostic experiments show $\text{grad\_norm\_central} = 0.0000$, meaning $c$ and $W_t$ receive zero gradient and freeze at a suboptimal value.

The root cause is \textit{tanh saturation}. The temperature scalar is $\tau = 1 + \epsilon \cdot \tanh(t_{\text{scalar}})$. When $|t_{\text{scalar}}|$ grows large, $\tanh(t_{\text{scalar}}) \approx \pm 1$ and its derivative vanishes:
\begin{equation}
    \frac{\partial \tau}{\partial t_{\text{scalar}}} = \epsilon \cdot (1 - \tanh^2(t_{\text{scalar}})) \approx 0
\end{equation}
By the chain rule, $\partial \mathcal{L}/\partial c \approx 0$. This is a hidden failure: the forward pass runs correctly, but the backward pass is severed.

\section{Experiments}

\subsection{Setup}

CPU sandbox: Docker on WSL2, Debian 12, Python 3.11.2, PyTorch 2.13.0+cpu, AMD Ryzen 9 8945HX 32 cores, $D \leq 128$. Results averaged over 5 seeds.

\subsection{Main Results}

Table~\ref{tab:main_results} summarizes the findings. Adaptive-EMA achieves the lowest MSE on v22 (0.2972), but gate-style wins at medium and large corpus scales. The crossover occurs at 16--32 training samples.

\begin{table}[ht]
\centering
\caption{Main results: four-mode comparison (v22) and corpus-scale boundary (v29).}
\label{tab:main_results}
\begin{tabular}{lcc}
\toprule
\textbf{Setting} & \textbf{Best Mode} & \textbf{MSE} \\
\midrule
v22 cross-architecture & Adaptive-EMA & 0.2972 \\
v29 tiny (8/4/4) & Adaptive-EMA & 0.7035 \\
v29 small (18/6/6) & Adaptive-EMA & 0.6984 \\
v29 medium (32/12/12) & Gate & 0.6971 \\
v29 large (64/24/24) & Gate & 0.7029 \\
\bottomrule
\end{tabular}
\end{table}

\textbf{Held-out correction.} The table above reports in-sample errors. Re-running the four hub modes against a held-out split of the same distribution (\texttt{run\_v44\_hub\_modes\_audit.py}, 5 seeds, paired per seed) changes the ranking: router-norm becomes the best mode out of sample (held-out $0.3342$, ahead of broadcast on $5/5$ paired seeds), while adaptive-ema---the in-sample winner here---collapses to $+0.3\%$ over broadcast and loses on $2/5$ seeds, with the largest generalization gap of the four ($+25.0\%$). An independent re-run on real corpora at a different scale (\texttt{run\_v25\_corpus\_central.py}, which already computed held-out scores but never summarized them) reproduces the ordering: router-norm best on held-out ($+4.6\%$ over broadcast), adaptive-ema second ($+2.0\%$), gate negative ($-5.1\%$). The in-sample ranking of this table should therefore be read as a training-fit ordering, not a capability ranking.

\textbf{The corpus-scale boundary does not survive the same test.} Sweeping only the two endpoint scales (\texttt{run\_v45\_scale\_boundary\_audit.py}, small $=v26$ configuration and large $=v29$ configuration; 5 seeds; train and held-out), the boundary is present in-sample (adaptive-ema $\rightarrow$ gate) but the held-out medians give adaptive-ema at both scales, i.e.\ no boundary. Worse, at the large scale the two readouts disagree with each other: the unpaired median ranks adaptive-ema first while the paired per-seed comparison ranks gate first (adaptive-ema wins only $1/5$), the disagreement being driven by a single seed where adaptive-ema reaches $0.767$ against gate's $0.543$. We therefore report the boundary as \emph{not established out of sample} rather than confirmed: two of the four scale levels were not re-run (the crossover level itself was not), and the small-scale held-out margin is $+0.5\%$ on $3/5$ seeds, i.e.\ within noise. Establishing or breaking the boundary requires the medium tier and a paired protocol; it is not settled by the present data.

\textbf{Diagnosis:} The regression is the tanh-saturation failure mode diagnosed in Section~\ref{sec:diagnosis}: the hub stops updating ($\text{grad\_norm\_central} = 0.0000$), so gate-style silently degrades to a frozen, suboptimal temperature.

\textbf{Temperature noise hypothesis:} We test whether temperature fluctuations cause gate-style regression. Across three scales, $\tau$ reaches $[0.9, 1.1]$ with standard deviation $0.09$--$0.10$, which does not increase with scale. We \textbf{reject} the hypothesis that temperature noise accumulation causes regression.

\section{Discussion}
\label{sec:discussion}

Our experiments suggest gate-style is not universally optimal; its effectiveness depends on data scale and training horizon. Based on 29 versions and 175+ experiments, we summarize three applicability boundaries:

\begin{itemize}
    \item \textbf{Optimal}: Medium corpus ($\geq 32$ samples), medium training steps (100--300).
    \item \textbf{Suboptimal}: Small corpus ($\leq 16$ samples) or long-horizon training; adaptive-EMA is more stable.
    \item \textbf{Failure mode}: Long training with large $|t_{\text{scalar}}|$ without gradient protection. Tanh saturation is hidden and may not surface until performance regresses.
\end{itemize}

The key insight is that \textit{how} a hub interacts with routing matters more than \textit{whether} it exists. Additive bias fails because it violates order preservation; temperature modulation preserves order by construction. This suggests a broader principle: when inserting coordination signals into discrete selection mechanisms, multiplicative modulation is safer than additive perturbation.

\section{Conclusion}

We introduced Gate-Style Central Coordination, where a global hub controls softmax temperature rather than perturbing router logits. This guarantees perfect degeneration, preserves logit ordering, and satisfies five design constraints (G1--G5) that other hub modes violate. Through 29 versions and 175+ experiments, we diagnosed a hidden failure mode---tanh saturation severs gradients in long training---and established a corpus-scale boundary: gate-style is optimal for $\geq 32$ samples, while adaptive-EMA dominates under extreme scale constraints.

The applicability boundaries of Section~\ref{sec:discussion}---optimal at medium corpus and training horizon, adaptive-EMA more stable under small corpora or long-horizon training, and the hidden tanh-saturation failure mode---delimit when gate-style should be deployed. The key insight is that \textit{how} a hub interacts with routing matters more than \textit{whether} it exists: multiplicative modulation is safer than additive perturbation when inserting coordination signals into discrete selection mechanisms.

\section*{Reproducibility Statement}

All scripts and evaluation code are publicly available. Core reproduction commands:

\begin{itemize}
    \item \texttt{python3 examples/run\_v22\_theory\_compare.py}
    \item \texttt{python3 examples/compare\_v9\_v30\_v31.py}
    \item \texttt{python3 examples/diagnose\_v29\_gate\_regression.py}
\end{itemize}

Environment: Docker on WSL2, Debian 12, Python 3.11.2, PyTorch 2.13.0+cpu, AMD Ryzen 9 8945HX 32 cores, $D \leq 128$, 5 seeds.

\begin{ack}
% 匿名投稿时留空，camera-ready 时填写
\end{ack}

\bibliographystyle{plainnat}
\bibliography{references}

\appendix
\section{Proofs}
\label{app:proofs}

Throughout, $\epsilon \in (0,1)$, the router logits $z \in \mathbb{R}^M$ are arbitrary, and the no-hub baseline is the map $z \mapsto \softmax(z)$.

\medskip
\noindent\textbf{Theorem (Perfect Degeneration).} \emph{When $c = 0$, gate-style is exactly equivalent to the no-hub baseline.}

\begin{proof}
If $c = 0$, then $t_{\text{scalar}} = W_t^\top c = 0$ for every $W_t$, and $\tanh(0) = 0$, so $\tau = 1 + \epsilon \cdot 0 = 1$. The gate-style logits are $z_m^{(\text{gate})} = z_m / \tau = z_m$, hence $\softmax(z^{(\text{gate})}) = \softmax(z)$ pointwise for every logit vector. The hub leaves the routing distribution---and therefore the fused output---identical to the no-hub system.
\end{proof}

\medskip
\noindent\textbf{Theorem (Order Preservation).} \emph{Gate-style preserves logit ordering; broadcast may reverse it.}

\begin{proof}
Since $\tanh \in (-1, 1)$, we have $\tau = 1 + \epsilon \tanh(t_{\text{scalar}}) \in (1 - \epsilon,\, 1 + \epsilon) \subset (0, \infty)$. Dividing by one common positive scalar preserves order: if $z_1 > z_2$ then $z_1 / \tau > z_2 / \tau$, so the softmax ranking is unchanged.

For broadcast, let $M \geq 2$, pick any $c \neq 0$ and any $\delta > 0$. Because $U$ is unconstrained, choose rows $U_1, U_2$ with $(U_1 - U_2)\,c = -(z_1 - z_2) - \delta$. Then $z_1 + (Uc)_1 < z_2 + (Uc)_2$ even though $z_1 > z_2$: the ranking of the two experts is reversed. Broadcast therefore does not preserve ordering in general.
\end{proof}

\medskip
\noindent\textbf{Theorem (Unification).} \emph{Among the four hub modes of \S\ref{sec:modes}, gate-style is the only one satisfying all five constraints (G1)--(G5).}

\begin{proof}
We show each competing mode violates at least one constraint, then verify that gate-style satisfies all five.

\emph{Broadcast violates (G2)} by the counterexample in the proof of Order Preservation.

\emph{Router-Norm violates (G1).} With $c = 0$ the mode applies $z \mapsto \operatorname{LayerNorm}(z)$, which subtracts the mean and rescales by $\gamma / \sigma(z)$. Translation invariance of softmax makes the mean subtraction harmless, but the rescaling generally differs from the identity. For example, with $M = 2$, $z = (1, 0)$, and default $\gamma = 1$: $\mu = 1/2$, $\sigma = 1/2$, so $\operatorname{LayerNorm}(z) = (1, -1)$, while $\softmax(1, 0) \approx (0.731, 0.269)$ but $\softmax(1, -1) \approx (0.881, 0.119)$. With a zero hub the routing distribution is not the no-hub baseline, so degeneration is not exact.

\emph{Adaptive-EMA violates (G1).} At $c = 0$ the hub-controlled decay is $\sigmoid(\alpha \cdot \|0\|) = 1/2$, which in general differs from the no-hub baseline's decay. Moreover, the EMA state persists across steps and is updated under hub control, so the routing distribution at step $t$ depends on the hub through the state it has shaped; disabling the hub does not reproduce the baseline exactly. The same state carries content into the logits, so ordering can be reversed as in broadcast, violating (G2) as well.

\emph{Gate-style satisfies all five.} (G1) holds by Perfect Degeneration. (G2) holds by Order Preservation. (G3): $\tau$ moves monotonically away from $1$ as $|t_{\text{scalar}}|$ grows, so hub intensity sharpens or flattens the softmax deterministically. (G4): $\tau \in (1 - \epsilon,\, 1 + \epsilon)$ is bounded, so the hub's influence on the logits is bounded by construction. (G5): the hub's parameters are $W_t \in \mathbb{R}^{D_{\text{shared}}}$, i.e., $O(D_{\text{shared}})$, negligible relative to the experts. Hence gate-style is the only one of the four modes satisfying all five constraints.
\end{proof}

\newpage
\input{checklist.tex}

\end{document}
